% Synthetic DCE - combined methods & results report
% Build:  bash scripts/build_docs.sh   (or: pdflatex REPORT.tex, twice)
\documentclass[11pt]{article}
\usepackage[margin=1in]{geometry}
\usepackage{amsmath,amssymb,booktabs,microtype,longtable,array}
\usepackage[table]{xcolor}
\usepackage[colorlinks=true,linkcolor=blue!45!black,urlcolor=blue!45!black,
            citecolor=blue!45!black]{hyperref}

\newcommand{\yhat}{\hat{y}}
\newcommand{\roi}{\mathcal{M}}
\newcommand{\Var}{\operatorname{Var}}
\newcommand{\Ex}{\mathbb{E}}
\newcommand{\1}{\mathbb{1}}
\newcommand{\flag}[1]{\textcolor{red!60!black}{\textbf{#1}}}

\title{Synthetic DCE from Non-Contrast Prostate MRI\\[2pt]
       \large Methods and Results --- UCSF cohort, Tier-1 static track}
\author{}
\date{6 August 2026}

\begin{document}
\maketitle
\vspace{-2.5em}

\begin{abstract}
\noindent
We predict a post-contrast DCE volume from three non-contrast sequences (T2w, DWI, ADC)
across 3{,}667 UCSF prostate exams, motivated by gadolinium avoidance in
GFR-contraindicated patients and by cross-protocol harmonization. A 3D conditional GAN
attains within-gland co-localization $r=0.482$ against $-0.043$ for the strongest trivial
baseline, establishing that the model predicts enhancement structure rather than passing
anatomy through. The central result is a decomposition: within-gland \emph{structure} is
predictable, per-patient enhancement \emph{amplitude} is not --- and amplitude carries
$51\%$ of the signal. That deficit is then largely repaired by supplying the pre-contrast
phase as a fourth input channel, which requires no gadolinium: brightness tracking rises from
$r=0.334$ to $0.76$ and the model clears a pre-contrast passthrough control by $+0.267$
\texttt{roi\_pearson}. Amplitude is therefore unidentifiable from bpMRI alone rather than
unidentifiable in principle. Methods and the measurements that justify them are presented
together throughout.
\end{abstract}

\tableofcontents
\vspace{1em}

\section{Task and notation}

Let $x=(x^{\text{T2w}},x^{\text{DWI}},x^{\text{ADC}})\in\mathbb{R}^{3\times D\times H\times W}$
be the stacked non-contrast input and $y\in[-1,1]^{1\times D\times H\times W}$ the target DCE
phase. We learn $G_\theta:x\mapsto\yhat$, $\yhat\in[-1,1]$ via $\tanh$. $\roi_i$ is the
prostate mask of case $i$, occupying $\approx5\%$ of the crop; $\Omega$ is the full volume.
All intensities live in $[-1,1]$, so SSIM and PSNR use data range $L=2$.

\section{Data}

\subsection{Source and cohort construction}

Single-center UCSF, delivered across two secured mounts: anatomy and masks on DS2
(5{,}962 patients), 4D DCE with per-phase timing on DS3. All volumes arrive pre-registered
into T2w space. Filters are applied at two points --- the enhancement criterion at staging,
so the tree is correct by construction, and two load-time flags.

\begin{center}
\begin{tabular}{llr}
\toprule
Stage & Criterion & Remaining \\
\midrule
Patients with DCE & --- & 5{,}439 \\
Successfully staged & inputs resolvable, mask matches DCE grid & 4{,}975 \\
Enhancement QC (at staging) & $\max_t \text{enh}(t)\ge1.5$ & \textbf{4{,}246} \\
DWI b-value filter & DWI resolves to $b\ge600$ & 3{,}684 \\
Enhancement QC recorded & --- & \textbf{3{,}667} \\
\midrule
\multicolumn{2}{l}{Patient-level split (train / val / test)} & 2805 / 312 / 550 \\
\bottomrule
\end{tabular}
\end{center}

\paragraph{Two attrition steps are findings, not bookkeeping.}
Roughly $12$--$13\%$ of patients never enhance at \emph{any} timepoint. Enhancement QC keys
on the maximum over the whole series rather than the selected phase, because only the
maximum separates a genuinely non-enhancing study (drop) from a mis-selected phase
(recoverable). \flag{Caveat:} this is a mask-mean ratio, so a mis-registered mask depresses
it exactly like a failed injection; we now capture per-case registration quality but have
not yet correlated the two, so the ``failed injection'' reading is not established.

Separately, $13\%$ of DWI resolves to $b{=}50$, which carries almost no diffusion weighting
and reads T2-like --- a \emph{different contrast} occupying the same input channel as the
$b{=}1000$ the majority receive. Excluding it costs $13\%$ of the data and is a clean
one-flag ablation. A further 43 cases use \emph{computed} (\texttt{\_synth\_}) high-$b$
images rather than acquired ones.

\subsection{Target phase: selected by time, not index}

Enhancement relative to pre-contrast is
\begin{equation}
\text{enh}(t)=\frac{\Ex_{v\in\roi}[S_v(t)]}{\Ex_{v\in\roi}[S_v(0)]}.
\end{equation}
\textbf{Measured:} flat at $\approx1.0$ until $t\approx45$\,s; wash-in $55$--$90$\,s
($1.12\to2.00$); then a \textbf{plateau} of $2.1$--$2.5$ from $\approx90$\,s out past
$300$\,s. There is no sharp peak, so $\arg\max_t\text{enh}(t)$ lands on plateau noise at
$\approx280$\,s, drifting into washout.

Protocols are heterogeneous --- $26$--$70$ phases at $1.7$--$13$\,s cadence --- so a fixed
\emph{index} is a different physiological moment per patient. We therefore select by time,
\begin{equation}
j^\star=\arg\min_j\bigl|t_j-t_{\text{target}}\bigr|,\qquad t_{\text{target}}=120\,\text{s},
\end{equation}
inside the stable plateau, past the timing-sensitive wash-in, at $\approx92\%$ of maximum.
Realized selection times cluster tightly at $104$--$125$\,s.

\paragraph{Two protocol complications, handled explicitly.}
\emph{Interleaved series} (27 cases): some studies stack $k$ volumes per timepoint (Dixon
water/fat), so $n_{\text{phases}}=k\cdot n_{\text{timepoints}}$ and a plain index alternates
between sub-series, one of which barely enhances. Detected via repeated timestamps; the
phase is re-picked within the timepoint, each sub-series scored against its own $t{=}0$.
\emph{Corrupt timing} (155 cases): some timing files carry impossible trailing values
(e.g.\ 67989\,s). Only the clean monotonic prefix is trusted, and where the time-selected
phase still looks weak the phase is re-selected from the enhancement curve at plateau onset
(first phase $\ge0.9\max$).

\subsection{Staging}

Each 4D series is $\approx1$\,GB; training directly against them re-reads the full 4D per
sample per epoch and depends on a mount that has historically dropped on GPU nodes. Staging
reads each 4D \emph{once}, writes the selected phase as 3D, and co-locates every modality in
a flat per-patient directory with a metadata sidecar. The staged target is verified
bit-identical to the raw-4D path. Current tree: 463\,GB over 4{,}251 cases.

\section{Preprocessing}

\subsection{Geometry: a fixed field of view, not a fixed pixel count}

\textbf{Measurement first.} UCSF in-plane resolution is not uniform --- matrices of
$256/512/672/704/736/1024$ at $0.176$--$0.781$\,mm. A fixed \emph{pixel} crop of 256
therefore covers a wildly variable physical extent (full cohort, $n=4246$):

\begin{center}
\begin{tabular}{lrr}
\toprule
256-px crop covers & Cases & Share \\
\midrule
$<80$\,mm & 734 & $17\%$ \\
$80$--$110$\,mm (intended) & 3346 & $79\%$ \\
$>110$\,mm & 164 & $4\%$ \\
\bottomrule
\end{tabular}
\end{center}

An $8\times$ spread in anatomical scale ($40$--$320$\,mm), with 172 cases cropped to
$45$\,mm --- narrower than a prostate, so the crop sits \emph{inside} the gland, the mask
fills the frame, and ROI metrics degenerate to whole-image metrics. The through-plane axis
has the same defect in miniature: $97.8\%$ of cases are $3.00$\,mm slices, but 44 are
$0.70$--$0.80$\,mm, giving those a $25.6$\,mm slab where everyone else gets $96$\,mm.

\textbf{Method.} Resample all modalities to a common physical grid before cropping,
$(0.47,0.47,3.0)$\,mm for 3D, then center-crop to $32\times192\times192$: a fixed
$90\times90\times96$\,mm field of view for every patient. Validated as a near-no-op for the
dominant $512@0.352$\,mm group and the $97.8\%$ at $3.00$\,mm, and T2w/DCE grids are
identical in all 4{,}246 cases so the iso anchor introduces no crop-center shift.

\subsection{Intensity normalization}

Single center, so there is no cross-scanner distribution to align and no Nyul landmark
fitting is used. DCE is robust-scaled about the body-tissue median $m$ with spread $s$:
\begin{equation}
y\;\longmapsto\;\operatorname{clip}\!\Bigl(\frac{y-m}{k\,s},-1,1\Bigr),\qquad k=1 .
\end{equation}
\textbf{Justification.} Within-prostate target dispersion measured over 12 staged cases:
$k{=}3\Rightarrow0.073$, $k{=}2\Rightarrow0.109$, $k{=}1\Rightarrow\mathbf{0.218}$. Larger
$k$ flattens the target until \texttt{roi\_pearson} and \texttt{ssim\_roi} become noise on a
near-uniform field. T2w uses percentile normalization, ADC a fixed physical clip (preserving
absolute mm$^2$/s), DWI a per-image foreground $z$-score.

\subsection{Split}

Patient-level, seeded, $15\%$ held out; validation carved from the training split only.
\flag{Note:} \texttt{-{}-seed} currently controls both the split and training
initialization, so multi-seed repeats vary both unless a separate split seed is added.
ID prefixes were tested as a possible batch marker and rejected (\texttt{8ii*} spans 512,
704, 736 and 1024 matrices alike).

\section{Models and objective}

\paragraph{Families.} 3D conditional GAN (pix2pix-style U-Net generator with full skips,
hinge loss, projection discriminator); 3D latent flow matching and DDPM sharing a frozen
first stage --- either a trained \texttt{AutoencoderKL3D} or the frozen MedVAE foundation
VAE. Every trainer returns $(model, \texttt{gen})$ with $\texttt{gen}(x)\to\yhat$, so
evaluation is model-agnostic. A 2D track (pix2pix, pixel-space CFM, MedVAE-latent flow)
turns $\approx3.2$k volumes into $\approx52$k slices at $\approx6\times$ lower cost per step.

\paragraph{Objective.} With $\phi$ a frozen 3D MedicalNet feature extractor,
\begin{equation}
\mathcal{L}=\lambda_1\mathcal{L}_1^{\text{ROI}}
+\lambda_s\bigl(1-\mathrm{SSIM}_{3D}\bigr)
+\lambda_s\bigl(1-\mathrm{SSIM}_{\roi}\bigr)
+\lambda_p\bigl\|\phi(\yhat)-\phi(y)\bigr\|_2^2 ,
\end{equation}
\begin{equation}
\mathcal{L}_1^{\text{ROI}}=\frac{\sum_{v\in\Omega}w_v|\yhat_v-y_v|}{\sum_{v\in\Omega}w_v},
\qquad w_v=1+(\omega-1)\1[v\in\roi],\qquad \omega=10 .
\end{equation}
Without ROI weighting the loss is dominated by background and padding.
\flag{Asymmetry worth stating:} the latent models train on latent-space MSE with no clean
mapping to the mask, so ROI emphasis reaches them only through the first stage's
reconstruction. Any GAN-versus-flow comparison inherits this.

\paragraph{Flow matching.} On the rectified path with $t{=}0$ at the data and $t{=}1$ at the
source,
\begin{equation}
z_t=\bigl(1-(1-\sigma_{\min})t\bigr)z_0+t\epsilon,\qquad
v^\star=\epsilon-(1-\sigma_{\min})z_0,\qquad \epsilon\sim\mathcal{N}(0,I),
\end{equation}
$\sigma_{\min}=0$, $t\sim\mathcal{U}(0,1)$ per example. The clean-latent identity
$\hat z_0=z_t-t\,v_\theta(z_t,t)$ holds exactly for any $\sigma_{\min}$ and is what the
optional image-space anchor decodes and supervises.

\section{Metrics}

\paragraph{Everything ROI is per-case, then averaged.} Pooling masked voxels across the
batch mixes between-patient brightness into every statistic. On synthetic data where
per-patient brightness is predicted perfectly and within-gland structure is pure noise,
pooling reports $r=+0.977$ where the truth is $+0.020$, and the value climbs with batch size
($0.66$ at $B{=}2$, $0.98$ at $B{=}8$). This was a live bug; it is now pinned by assertion.

\begin{align}
r_i&=\frac{\sum_{v\in\roi_i}(\yhat_v-\bar{\yhat})(y_v-\bar y)}
{\sqrt{\sum_{v\in\roi_i}(\yhat_v-\bar{\yhat})^2}\sqrt{\sum_{v\in\roi_i}(y_v-\bar y)^2}},
&\texttt{roi\_pearson}&=\tfrac1N\textstyle\sum_i r_i, \\[3pt]
\texttt{roi\_var\_ratio}&=\tfrac1N\textstyle\sum_i\min\!\Bigl(\tfrac{\Var_{v\in\roi_i}[\yhat]}{\Var_{v\in\roi_i}[y]+\varepsilon},5\Bigr),
&\texttt{roi\_grad\_ratio}&=\tfrac1N\textstyle\sum_i\min\!\Bigl(\tfrac{\Ex_{v\in\roi_i}\|\nabla\yhat\|}{\Ex_{v\in\roi_i}\|\nabla y\|+\varepsilon},5\Bigr),
\end{align}
both with target $1$ ($<1$ over-smoothed, $>1$ over-textured). The histogram distance uses
order statistics, $\texttt{roi\_w1}=\tfrac12|\roi_i|^{-1}\sum_k|\yhat_{(k)}-y_{(k)}|$, and a
label-free composite
\begin{equation}
\texttt{realism}=\operatorname{mean}\bigl\{\nu(\texttt{var\_ratio}),\,\nu(\texttt{grad\_ratio}),\,1-\min(\texttt{w1},1)\bigr\},
\quad \nu(x)=1-\min(|x-1|,1).
\end{equation}
\flag{The composite is a ranking heuristic, not a validated metric} --- equal weights, no
weighting study. Report components individually. \flag{The $\min(\cdot,5)$ caps bias both
ratios} toward 5 from above; they exist because a near-flat target denominator sent uncapped
values into the thousands.

\paragraph{Global SSIM/PSNR are $\approx99\%$ background} and near-meaningless here. The
\emph{gap} between global and ROI values is the diagnostic.

\section{Results}

\subsection{Model comparison}

3D models, best checkpoint, held-out test ($n=550$). VAL $\approx$ TEST throughout.

\begin{center}
\begin{tabular}{lccc}
\toprule
& 3D GAN & 3D flow (MedVAE) & GAN $+$ pre-contrast \\
\midrule
\texttt{roi\_pearson} $\uparrow$ & 0.482 & 0.369 & \textbf{0.527} \\
\texttt{ssim\_roi} $\uparrow$ & 0.397 & 0.276 & \textbf{0.449} \\
$\mathrm{MAE}_{\roi}$ $\downarrow$ & 0.236 & 0.269 & \textbf{0.200} \\
\texttt{roi\_var\_ratio} $\to1$ & 0.912 & 1.711 & \textbf{0.951} \\
\texttt{roi\_grad\_ratio} $\to1$ & 0.711 & \textbf{1.059} & 0.773 \\
\texttt{p75\_corr} $\uparrow$ & 0.357 & 0.351 & \textbf{0.743} \\
\texttt{realism} $\uparrow$ & 0.845 & 0.714 & \textbf{0.887} \\
FID $\downarrow$ & 88.9$^\dagger$ & 169.8$^\dagger$ & \textbf{84.6} \\
PZ / TZ \texttt{roi\_pearson} & 0.339 / 0.522 & 0.169 / 0.417 & \textbf{0.445} / \textbf{0.547} \\
\bottomrule
\end{tabular}
\end{center}
{\footnotesize $^\dagger$ measured on the \emph{last} checkpoint, not the best --- see
\S\ref{sec:corrections}.}

The GAN leads the flow on faithfulness across four independent setups (2D and 3D, pixel and
latent). The flow's one advantage is texture ($\texttt{grad\_ratio}=1.06$ vs $0.71$): it is
less smooth, which reads as more realistic while localizing worse.
\flag{Two confounds:} checkpoints are selected on \texttt{ssim\_roi}, which is
smoothness-biased and therefore favours the GAN's failure mode; and the flow's objective is
effectively ROI-unweighted (\S4). A selection-metric ablation has not been run.

\paragraph{Peripheral zone.} $r_{\text{PZ}}=0.339$ against $r_{\text{TZ}}=0.522$ for the
base GAN. PZ is where most clinically significant cancer arises and where DCE drives the
PI-RADS upgrade rule, so this is the clinically weakest region. Pre-contrast conditioning
lifts PZ disproportionately ($+31\%$ against TZ's $+5\%$).

\subsection{The identifiability decomposition}
\label{sec:decomp}

Three predictors requiring no learning, with $\mu_i=\Ex_{v\in\roi_i}[y_v]$,
$\bar\mu=\Ex_i[\mu_i]$:
$\yhat^{\text{const}}_v=\bar\mu$;
$\yhat^{\text{level}}_v=\mu_i$ (an \emph{oracle} knowing correct per-case brightness with no
structure); $\yhat^{\text{t2w}}_v=x^{\text{T2w}}_v$.

\begin{center}
\begin{tabular}{lccl}
\toprule
Predictor & $\mathrm{MAE}_{\roi}$ & \texttt{roi\_pearson} & Knows \\
\midrule
\texttt{const} & 0.2510 & --- & nothing \\
\texttt{t2w copy} & 0.3372 & $\mathbf{-0.043}$ & identity baseline \\
\textbf{MODEL} & 0.2355 & $\mathbf{+0.482}$ & trained generator \\
\texttt{level} (oracle) & \textbf{0.1851} & --- & brightness only, no structure \\
\bottomrule
\end{tabular}
\end{center}

Total ROI variance decomposes exactly, by the law of total variance:
\begin{equation}
\Var_{i,v}[y]=\underbrace{\Var_i[\mu_i]}_{\text{amplitude}}+\underbrace{\Ex_i[\sigma_i^2]}_{\text{structure}},
\qquad \sigma_i^2=\Var_{v\in\roi_i}[y_v].
\end{equation}
Measured: $\operatorname{sd}_i[\mu_i]=0.231$, $\Ex_i[\sigma_i]=0.226$, giving an amplitude
share of $0.231^2/(0.231^2+0.226^2)=\mathbf{0.51}$.
\flag{This is an upper bound:} we substituted $(\Ex[\sigma])^2$ for $\Ex[\sigma^2]$, and
Jensen makes the true structure term at least that large.

\paragraph{Reading it.}
\emph{Structure is learned} --- T2w-copy scores $\approx0$, so $r=0.482$ is genuine
localization, not laundered anatomy. \emph{Amplitude is not} --- the model tracks per-patient
brightness at $\operatorname{corr}_i(\mu_i,\hat\mu_i)=0.334$ and compresses between-case
spread $2.4\times$ ($0.231\to0.098$). Because amplitude is half the variance, both models
lose to the oracle-brightness baseline in MAE, and the flow is worse than a \emph{constant}
predictor ($-0.018$) while retaining real localization ($r=0.369$).

\paragraph{Why this is the task, not the model.} Per-case amplitude depends on gadolinium
dose, injection rate, cardiac output, renal clearance and receiver gain --- none visible in
T2w/DWI/ADC. We inspected the source DICOM headers for exactly these variables: they were
removed by de-identification (dose, rate, patient weight and field strength absent;
manufacturer overwritten with \texttt{"NA"}; only patient age survives). The information is
not merely unused --- it is absent.

\flag{Revised below.} \S6.3 shows this deficit is largely repaired once the pre-contrast
phase is available: the claim holds for bpMRI alone, not in general.

\paragraph{Why it may not matter clinically.} PI-RADS curve types are scale-invariant:
$\text{type}(\alpha C(t))=\text{type}(C(t))$ for any $\alpha>0$, since Type I (progressive),
II (plateau) and III (washout) are defined by post-peak \emph{shape}. The unidentifiable
component is orthogonal to the endpoint.

\subsection{Pre-contrast conditioning}

Supplying the pre-contrast phase (DCE phase 0, same acquisition and geometry) as a fourth
input channel, identical configuration otherwise:

\begin{center}
\begin{tabular}{lccr}
\toprule
& Base GAN & $+$ pre-contrast & $\Delta$ \\
\midrule
\texttt{roi\_pearson} & 0.482 & 0.527 & $+0.045$ ($\approx2.6\times$ noise) \\
\texttt{p75\_corr} (amplitude) & 0.357 & 0.743 & $+108\%$ \\
PZ \texttt{roi\_pearson} & 0.339 & 0.445 & $+31\%$ \\
$\mathrm{MAE}_{\roi}$ & 0.236 & 0.200 & $-15\%$ \\
Global SSIM / PSNR & 0.469 / 16.8 & 0.697 / 21.1 & $+49\%$ / $+4.3$\,dB \\
\bottomrule
\end{tabular}
\end{center}

\paragraph{Passthrough control.} Pre-contrast is a far stronger predictor of post-contrast
than T2w --- same acquisition, same receiver gain, post $\approx$ pre $+$ enhancement --- so
a model handed that channel could depress MAE simply by copying it. Adding
$\yhat_v=x^{\text{pre}}_v$ to the baseline panel settles it:

\begin{center}
\begin{tabular}{lccl}
\toprule
Predictor & $\mathrm{MAE}_{\roi}$ & \texttt{roi\_pearson} & Knows \\
\midrule
\texttt{const} & 0.2510 & --- & nothing \\
\texttt{t2w copy} & 0.3372 & $-0.043$ & identity baseline \\
\texttt{PRE copy} & 0.2855 & $+0.260$ & pre-contrast passthrough \\
\textbf{MODEL} & \textbf{0.2005} & $\mathbf{+0.527}$ & trained generator \\
\texttt{level} (oracle) & 0.1851 & --- & brightness only, no structure \\
\bottomrule
\end{tabular}
\end{center}

The model clears passthrough by $+0.267$ \texttt{roi\_pearson} ($\approx16\times$ the noise
floor) and $-0.085$ MAE. It is not copying: pre-contrast is used as a calibration reference
\emph{and} structure is added on top.

\paragraph{Amplitude is substantially recoverable --- a correction.}
\S\ref{sec:decomp} argued that per-case amplitude is unidentifiable. That was too strong: it
is unidentifiable \emph{from bpMRI alone}. Every stable amplitude readout moves sharply once
pre-contrast is available.

\begin{center}
\begin{tabular}{lccc}
\toprule
& Base GAN & $+$ pre-contrast & Target \\
\midrule
brightness tracking $\operatorname{corr}_i(\mu_i,\hat\mu_i)$ & 0.334 & \textbf{0.76} & 1.0 \\
between-case predicted sd & 0.098 & \textbf{0.189} & 0.231 \\
bias sd & 0.218 & \textbf{0.151} & --- \\
gap to oracle brightness (MAE) & $+0.0505$ & \textbf{$+0.0154$} & $\le0$ \\
improvement over \texttt{const} (MAE) & $+0.0154$ & \textbf{$+0.0504$} & $>0$ \\
\bottomrule
\end{tabular}
\end{center}

Spread compression falls from $2.4\times$ to $1.22\times$, and the model closes from $27\%$
worse than the oracle-brightness baseline to $8\%$ worse.

\paragraph{The two-component split survives; the sizes were wrong.} Amplitude decomposes
into an \emph{acquisition/baseline} term (receiver gain, the tissue baseline the gland
enhances from) and an \emph{injection/physiology} term (dose, rate, cardiac output, renal
clearance). Pre-contrast --- same series, same gain --- supplies the first, and it turns out
to dominate. The residual ($r=0.76$ rather than $1.0$; spread still $18\%$ compressed; MAE
still short of the oracle) is plausibly the second, which de-identification removed from the
headers and which no image contains. Clinically this costs nothing: a pre-contrast $T_1$
requires no gadolinium, so the premise of the work is intact.

\flag{Remaining caveat.} These are single runs. The large movements (passthrough margin,
brightness tracking, spread) are many multiples of the noise floor; the
$+0.045$ \texttt{roi\_pearson} gain over the base GAN is only $\approx2.6\times$ and would
benefit from seed repeats.

\section{Reliability}

\subsection{Pipeline corrections}
\label{sec:corrections}

Every defect below was \emph{silent} --- no crash, no warning, plausible-looking output.
Grouped by remediation cost: measurement defects can be re-scored from checkpoints, training
defects need retraining, data defects need re-staging.

\begin{center}
\small
\begin{tabular}{p{3.4cm}p{7.2cm}p{3.4cm}}
\toprule
Defect & Effect & Evidence \\
\midrule
Pooled ROI statistics & cross-patient brightness masqueraded as within-gland correlation & $r=+0.977$ on pure noise where truth is $+0.020$ \\
Transposed convolutions & period-2 checkerboard stamped into every prediction & lag-1 autocorr.\ $-0.60$ vs $+0.81$ real; FID $175\to89$ \\
Pixel-count crop & $8\times$ range of anatomical scale & $40$--$320$\,mm across cohort \\
Final-epoch evaluation & 3D \texttt{metrics.json} scored a degraded model & \texttt{roi\_pearson} $0.419\to0.482$ \\
Corrupt-timing fallback & 155 cases trained against the \emph{pre-contrast} volume & re-selected at plateau onset \\
Re-stage without QC & 729 pruned non-enhancers silently resurrected & a run trained on 4{,}975 \\
Harmonizer leakage & Nyul fit on the full cohort incl.\ test & never triggered by our configs \\
\bottomrule
\end{tabular}
\end{center}

The checkerboard detector is the lag-1 autocorrelation of the first difference,
$\rho=\operatorname{corr}(d_k,d_{k+1})$ with $d_k=y_{k+1}-y_k$, which tends to $-1$ under a
pure checkerboard and $-0.5$ for white noise. Predictions scored $\mathbf{-0.598}$ ---
alternating in sign at pixel scale \emph{more strongly than white noise} --- against
$+0.813$ for real DCE, with $5\times$ the target's high-frequency power. Post-fix: $+0.278$
and HF power at the target's level.

A 28-assertion preflight (\texttt{python -m tier1\_static.selfcheck}) pins each of these.
The remaining leakage surface was audited clean: train/test disjoint and exhaustive; val
carved from train only; \texttt{test\_loader} never used during training; latent scaling
computed from train only; 2D splits patients before slicing; evaluation and validation share
clamping and metric code.

\subsection{Run-to-run variance}
\label{sec:variance}

Two \emph{configuration-identical} runs at the same seed differ materially. Seeding covers
Python/NumPy/Torch but not cuDNN autotuning or non-deterministic \texttt{atomicAdd} in 3D
convolution backward; adversarial dynamics amplify the divergence; and best-checkpoint
selection converts a bumpy validation curve into a noisy discrete choice of epoch.

\begin{center}
\begin{tabular}{lc}
\toprule
Metric & $\Delta$ between identical runs \\
\midrule
\texttt{roi\_pearson} & 0.017 \\
\texttt{ssim\_roi} & 0.008 \\
\texttt{roi\_var\_ratio} & 0.055 \\
\texttt{p75\_corr} & \textbf{0.276} \\
\bottomrule
\end{tabular}
\end{center}

\paragraph{Rule.} Differences below $\approx0.02$ \texttt{roi\_pearson} are unresolved at
$n=1$. \emph{Safe:} GAN $>$ flow ($0.113$, $\approx7\times$); model $\gg$ t2w
($\approx30\times$); PZ $<$ TZ (paired within run); best $>$ last checkpoint
($\approx4\times$). \emph{Not safe:} any \texttt{p75\_corr} comparison at $n=1$; the
iso-versus-no-iso A/B ($0.020$ --- inside noise, so iso is justified on correctness grounds
only). This is one replicate pair, bounding the magnitude but not a standard deviation;
$\ge3$ seeds are required.

\paragraph{Why \texttt{p75\_corr} is $16\times$ worse.} \texttt{roi\_pearson} averages $N=550$
per-case values, suppressing variance by $\sqrt N$. \texttt{p75\_corr} is a single
cross-patient correlation measuring the \emph{amplitude} dimension --- precisely what nothing
in the loss constrains. \textbf{Reproducibility tracks identifiability}: the well-determined
component reproduces, the ill-determined one drifts. The instability is a measurement of
which part of the task is ill-posed, not merely a nuisance.

\section{Tier-2: kinetics}

The multi-timepoint target is the per-voxel enhancement curve. Unlike generic video world
models, the governing dynamics are \emph{known} --- the Tofts model,
\begin{equation}
\frac{dC_t}{dt}=K^{\text{trans}}C_p(t)-k_{ep}C_t(t),\qquad k_{ep}=K^{\text{trans}}/v_e ,
\end{equation}
with $C_p$ the arterial input function. A continuous-time latent world model integrates a
learned velocity field in physical time,
\begin{equation}
\frac{dz_t}{dt}=v_\theta(z_t,t;z_0,c)\;\Longrightarrow\;z_T=z_0+\int_0^T v_\theta\,dt ,
\end{equation}
with $z_0$ the encoded pre-contrast state and $c$ a conditioning signal. The proposed
contribution is to \emph{constrain} $v_\theta$ to the Tofts family with parameter maps
predicted from bpMRI rather than learn it free-form --- a stronger inductive bias with
guaranteed physically plausible curves. This directly addresses the limitation such methods
cite, that physics-informed ODE approaches ``require explicit prior knowledge of the true
governing equations''; for DCE we have it.

\flag{Two obstacles to state up front.} Tofts governs \emph{concentration}; we have
normalized signal, and converting requires a baseline $T_1$ map this cohort lacks, so
recovered parameters are semi-quantitative at best. And an arterial input function is
required, which a $90$\,mm prostate-centred crop may not contain.

\paragraph{Goal state.} Acquisition duration (measured, $n=1367$): $p_5=222$\,s,
$p_{50}=300$\,s, $p_{95}=380$\,s. Coverage: $99.0\%$ at $180$\,s, $\mathbf{91.9\%}$ at
$240$\,s, $49.7\%$ at $300$\,s. We anchor at $240$\,s --- past plateau onset by
$\approx150$\,s, inside the window where Types II and III separate, retaining $92\%$ of the
cohort. The $8\%$ falling short run the same protocol at faster cadence (median 45 phases
versus 36), not a different scanner, so excluding them from a curve-type readout introduces
no protocol confound.

\section{Paper partitions and venues}

\begin{center}
\renewcommand{\arraystretch}{1.2}
\small
\begin{tabular}{p{0.5cm}p{4.2cm}p{4.6cm}p{3.6cm}}
\toprule
& Claim & Still needs & Venue \\
\midrule
1 & Single-phase synthetic DCE is clinically usable & labels; reader study; multi-seed & Radiology / Lancet Dig.\ Health \\
2 & Standard synthesis metrics do not measure case-specific fidelity & within-model sweep; $\ge3$ seeds; \emph{second domain} & ICLR (tight) $\to$ ICML / NeurIPS \\
3 & A physics-constrained velocity field beats a free-form one & Tofts parameterization; interpolation ablation & NeurIPS workshop $\to$ MICCAI \\
\bottomrule
\end{tabular}
\end{center}

Track 2 is the strongest near-term result and is \emph{label-free}: the decomposition of
\S\ref{sec:decomp}, the reproducibility-tracks-identifiability observation, and now a
constructive counterpart --- the missing component is recovered by naming the right input,
which turns ``here is something that cannot be predicted'' into ``here is which input carries
which component of the target,'' a stronger and more useful claim. Its framing
must avoid the perception--distortion tradeoff, which is established; the contribution is a
\emph{diagnostic for whether a conditional model used its conditioning at all}. A single
medical domain reads as an application paper, so a second domain is the gating requirement.

Track 3's novelty is not porting a world model to a new organ --- that is a MICCAI paper ---
but that we know the governing equation.

\paragraph{On the Tier-1 $\to$ Tier-2 gate.} The project's onboarding sets a hard gate at
PI-RADS $\kappa\ge0.6$. That is currently \emph{unmeasurable} rather than failed: the overall
PI-RADS category is essentially absent from the report export ($92/10{,}538$ once a
boilerplate footer URL is excluded), though the binary DCE component appears in
$\approx6{,}400$ reports. Two developments argue for revisiting it: the scarcity rationale
has evaporated (UCSF supplies $\approx8\times$ the planned Tier-2 pool at $\approx10\times$
the temporal density), and curve types are scale-invariant, so the one component we cannot
predict is orthogonal to the endpoint. Gating on \emph{DCE-positivity} $\kappa$ is the
defensible substitute.

\section{Anticipated technical questions}

\paragraph{Is $r=0.482$ inflated by zonal contrast rather than focal enhancement?}
A model learning only ``TZ is brighter than PZ'' would score well without detecting anything.
The data partly answers it: if the between-zone term were doing the work, whole-gland $r$
should greatly exceed both within-zone values. It does not --- $0.482$ sits \emph{between}
$r_{\text{PZ}}=0.339$ and $r_{\text{TZ}}=0.522$, and both within-zone values are solidly
positive. This does \emph{not} establish lesion-level detection, which needs lesion labels.

\paragraph{You report a mean over 550 cases with no interval.}
Correct, and it should be added --- the per-case $r_i$ distribution and a bootstrap CI are
computable from data in hand. The $\pm0.017$ figure bounds \emph{run} variability, not
\emph{case} variability; the two must not be conflated.

\paragraph{Why Pearson? The signal--concentration relationship is nonlinear.}
A real limitation. MR signal relates to concentration through $T_1$ relaxation and the
spoiled gradient-echo signal equation. Spearman $\rho$ is the rank-based alternative and is
trivial to add; reporting both would be more defensible.

\paragraph{Does the variance decomposition assume independence?}
No --- the law of total variance is an identity. The only caveat is the
$(\Ex[\sigma])^2$-for-$\Ex[\sigma^2]$ substitution, making $51\%$ an upper bound.

\paragraph{Is ``the GAN wins'' a checkpoint-selection artifact?}
Partly. \texttt{ssim\_roi} is smoothness-biased and the GAN is the smoother model. Against
that: the GAN also leads on \texttt{roi\_pearson}, which is not smoothness-biased, and by
$\approx7\times$ the noise. A selection-metric ablation is the clean answer and is not done.

\paragraph{Were the 3D runs converged? 40 epochs versus 100 for 2D.}
A wall-clock choice ($\approx670$\,s/epoch), not a convergence criterion, and convergence
was not demonstrated. Against it: the GAN \emph{degrades} late (discriminator wins,
$d\to0.03$), and the selected checkpoint clearly beats the final epoch. Per-run validation
trajectories are now logged and will answer this directly.

\paragraph{Are the non-enhancers biology or registration failure?}
Not separated. A mis-registered mask depresses mask-mean enhancement exactly like a failed
injection. Per-case registration quality is now captured (cohort median $0.459$,
$p_{90}=1.694$) but has not been correlated against enhancement.

\paragraph{FID on 550 grayscale volumes with an ImageNet backbone?}
Weak. FID is biased at small $n$ (KID is the alternative) and Inception features are
natural-image. It is retained because it responded sharply and correctly to a known artifact
($175\to89$), but it should not carry a headline claim.

\paragraph{How were the masks produced?}
Automatic segmentations delivered with the cohort; none manually verified. Every ROI number
here is conditioned on them.

\paragraph{Is the effective test sample really 550?}
Probably smaller. Protocols cluster (generic, Siemens TWIST, GE DISCO) and the split is not
stratified by protocol. Stratifying --- now possible, since the protocol string is captured
--- would fix this and double as the multi-vendor generalization result.

\section{Reproducing this}

\paragraph{Cohort.} Stage once with the QC filter inline --- applying it at staging rather
than pruning afterwards is what makes the tree correct by construction:
\begin{verbatim}
MIN_ENH=1.5 sbatch --array=0-3 scripts/stage_ucsf.slurm      # -> 4,246
python -m tier1_static.data.stage_ucsf --report              # verify before training
python -m tier1_static.selfcheck                             # 28 assertions, CPU, seconds
\end{verbatim}
Expect median enhancement $\approx2.07$, $p_{10}=1.70$, 27 interleaved, 155 curve-reselected,
and no $t{=}0$ warning. Training logs must show \texttt{min\_b=600} and
\texttt{562 below b600, 17 without enhancement QC}.

\paragraph{Training.}
\begin{verbatim}
UCSF_ROOT=<staged> MODEL=gan OUTPUT_DIR=runs/<name> \
REFERENCE=iso ISO_SPACING="0.47 0.47 3.0" SPATIAL_SIZE="32 192 192" \
T2W_NORM=percentile DCE_NORM=robust DCE_ROBUST_K=1 \
EPOCHS=40 BATCH_SIZE=4 BASE_CH=32 ROI_WEIGHT=10 VAL_EVERY=2 SEED=0 \
EXTRA_ARGS="--dwi-min-bvalue 600 --require-qc" sbatch scripts/train.slurm
\end{verbatim}
The flags must match between training, evaluation and audit, or the checkpoint is scored on a
different distribution than it was trained on. \texttt{config.json} is written by both the 2D
and 3D paths for exactly this reason. Each run also emits
\texttt{val\_trajectory.jsonl}, one row per validation step with the full metric dict --- the
faithfulness-versus-realism trajectory of \S\ref{sec:variance} is read from there.

\paragraph{Baseline audit.} Same flags plus \texttt{-{}-audit-split test}:
\begin{verbatim}
python -m tier1_static.audit_baselines --model gan ... --audit-split test
\end{verbatim}

\paragraph{Multi-timepoint staging.} Anchors on a common time grid, dense through wash-in and
sparse across the plateau, since that is where curve types separate:
\begin{verbatim}
--anchor-times 45 60 75 90 150 180 240      # 0s and 120s already staged
\end{verbatim}
Seven additional phases cost $\approx0.76$\,TB over the existing 463\,GB, and one read of
each 4D rather than seven.

\section{Limitations}

2D and 3D ROI metrics are not comparable (the 2D dataset keeps only gland-containing slices,
so denominators differ). FID for the base models was measured on the last checkpoint. Every
result predating the audit is superseded. No claim is made that more data would close the
amplitude gap: an earlier conclusion to that effect was measured on a defective configuration
and has been retracted. The pre-contrast result awaits its passthrough control.

\end{document}
